%% ============================================================
%% 分片 20：原文第 7 节「Fluxes and $S$-Duality」（原文 glshort.tex 起于第 5685 行）
%% 由 tools/harvmac2tex.py 从 arXiv 源码（hep-th/0604151v3）自动生成。
%% 只译英文散文；公式、编号、\label、\cite 一律不改。
%% ============================================================

\subsection{紧化到二维} \label{sec:comptwop}

现在我们把理论紧化到二维。考虑 $M=\Sigma\times C$ 上规范群为 $G_{\rm{ad}}$ 的 ${\EUN}=4$ 超 Yang-Mills 理论。该理论的数据是 $G_{\rm{ad}}$-丛 $E\to\Sigma\times C$、$E$ 上的联络 $A$ 以及各种其他场。任取 $E$ 都决定一个类 $\xi(E)\in H^2(\Sigma\times C,\CZ)$。由于 $H^2(\Sigma\times C,\CZ)=H^2(\Sigma,\CZ)\oplus H^1(\Sigma,H^1(C,\CZ))\oplus H^2(C,\CZ)$，我们有 % [WARN] 公式 \eqn\eqref{eq:ytro} 无编号或重复
\begin{equation*}
\tag*{(7.18)}
\xi(E)\in H^2(\Sigma,\CZ)\oplus
H^1(\Sigma,H^1(C,\CZ))\oplus
H^2(C,\CZ).
\end{equation*}
相对于这一分解，我们写 $\xi(E)=\xi(E)^{2,0}+\xi(E)^{1,1}+\xi(E)^{0,2}$。

我们可以用低能有效 sigma 模型描述同一件事。这里我们考虑 $\Sigma\times C$ 上的低能数据，当它们限制到 $p\times C$（$p\in \Sigma$）上时，决定 $\MHAD$ 中的一个点。随着 $p$ 变化，我们得到映射 $\Phi:\Sigma\to \MHAD$。把 $\Phi$ 与 $C$ 上的恒等映射复合，得到映射 $\Phi\times 1:\Sigma\times C\to \MHAD\times C$。现在回忆泛丛 ${\EUE}_{\rm ad}$ 与类 $\zeta=\xi({\EUE}_{\rm {ad}})\in H^2( \MHAD\times C,\CZ)$。把它拉回到 $\Sigma\times C$，我们得到类 $(\Phi\times 1)^*(\zeta)\in H^2(\Sigma\times C,\CZ)$。梳理一遍定义即可看出 % 式 (7.19)
\setcounter{equation}{18}
\begin{equation}\label{eq:yrog}
\xi(E)=(\Phi\times
1)^*(\zeta).
\end{equation}
事实上，泛丛 ${\EUE}_{ad}$ 正是这样定义的，使得 $(\Phi\times 1)^*({\EUE}_{{\rm ad}})=E$，所以 $(\Phi\times 1)^*(\zeta)=(\Phi\times 1)^*(\xi(\EUE_{{\rm ad}}))=\xi((\Phi\times 1)^*(\EUE_{{\rm ad}}))=\xi(E)$。

\begin{yizhu}[物理图像]
按：式 (7.18) 按 Künneth 分解把 $\xi(E)$ 分成 $\Sigma$ 与 $C$ 两个方向上的分量，$\xi(E)^{p,q}$ 即 $\Sigma$ 方向 $p$ 次、$C$ 方向 $q$ 次的部分。其中 $\xi(E)^{0,2}={\bf m}_0$ 度量 $E$ 限制到 $C$ 上的拓扑类型，$\xi(E)^{2,0}={\bf a}_0$ 是沿 $\tilde S^1$ 方向的数据；式 (7.19) 说明它们全部是从泛丛的障碍类 $\zeta$ 拉回得到的，这是后文把 $\zeta$ 解释为平坦 $B$-场的基础。
\end{yizhu}

我们进一步取 $\Sigma=S^1\times \tilde S^1$，其中两个因子都是圆周。我们写 $W=\tilde S^1\times C$，于是 $\Sigma\times C=S^1\times W$。$\xi(E)$ 有熟悉的展开 $\xi(E)={\bf a}+{\bf m}$，其中 ${\bf a}\in H^1(W,\CZ)$，${\bf m}\in H^2(W,\CZ)$。但现在我们还可以进一步展开 $H^1(W,\CZ)=H^1(\tilde S^1,\CZ)\oplus H^1(C,\CZ)$，于是有 ${\bf a}={\bf a}_0+{\bf a}_1$，其中 % 式 (7.20)
\setcounter{equation}{19}
\begin{equation}\label{eq:mimot}
\begin{aligned}{\bf
a}_0&\in H^1(\tilde S^1,\CZ) =\CZ,\\ {\bf a}_1& \in
H^1(C,\CZ).
\end{aligned}
\end{equation}
类似地，我们可以展开 $H^2(W,\CZ)=H^1(\tilde S^1,H^1(C,\CZ))\oplus H^2(C,\CZ)$，于是有 ${\bf m}={\bf m}_0+{\bf m}_1$，其中 % 式 (7.21)
\setcounter{equation}{20}
\begin{equation}\label{eq:nimot}
\begin{aligned}{\bf
m}_0&\in H^2(C,\CZ)=\CZ \\
                     {\bf m}_1&\in
                     H^1(\tilde S^1,H^1(C,\CZ))=H^1(C,\CZ).
\end{aligned}
\end{equation}

与展开 \eqref{eq:ytro} 相对照，我们看到
% 式 (7.22)
\setcounter{equation}{21}
\begin{equation}\label{eq:zimot}
\xi(E)^{2,0}={\bf a}_0,~\xi(E)^{0,2}={\bf m}_0.
\end{equation}

我们可以类似地分解与 $\bf a$ 对偶的特征标 $\bf e$。显然有 ${\bf e}={\bf e}_0+{\bf e}_1$，其中 ${\bf e}_0$ 是 ${\bf a}_0\in \CZ$ 的特征标，${\bf e}_1$ 是 ${\bf a}_1\in H^1(C,\CZ)$ 的特征标。因此，利用 $\CZ$ 的自对偶与庞加莱对偶， % 式 (7.23)
\setcounter{equation}{22}
\begin{equation}\label{eq:tew}
\begin{aligned} {\bf e}_0 & \in
\CZ^\vee=\CZ\\ {\bf e}_1&\in H^1(C,\CZ)^\vee=H^1(C,\CZ).
\end{aligned}
\end{equation}
鉴于 \eqref{eq:coco}，$S$-对偶下的变换为
% 式 (7.24)
\setcounter{equation}{23}
\begin{equation}\label{eq:niko}
\begin{aligned}({\bf e}_0,{\bf m}_0)&\to (-{\bf m}_0,{\bf
e}_0)\\ ({\bf e}_1,{\bf m}_1)&\to (-{\bf m}_1,{\bf e}_1).
\end{aligned}
\end{equation}

我们的下一个目标是在靶为 $\MHAD$ 的有效二维 sigma 模型中解释 ${\bf e}_0$、${\bf e}_1$、${\bf m}_0$ 与 ${\bf m}_1$。我们将同时考虑闭黎曼面 $\Sigma$ 的情形与 $\Sigma$ 带边界的情形。后一情形将使我们能够理解 \eqref{eq:niko} 对几何 Langlands 纲领的含义。

\par\medskip\emph{${\bf m}_0$ 的解释}

最容易解释的是 ${\bf m}_0=\xi(E)^{0,2}$。我们 sigma 模型的靶空间 $\MH(G_{\rm ad},C)$ 不连通。它的各连通分支由 $G_{\rm ad}$-丛 $E\to C$ 的拓扑类型标记。而这正是 ${\bf m}_0$ 所度量的。

\par\medskip\emph{${\bf e}_0$ 的解释}

一般地，考虑映到某个流形 $X$ 的映射 $\Phi:\Sigma\to X$ 的 sigma 模型。平坦 $B$-场是元素 $b\in H^2(X,U(1))$。在 $\Sigma$ 无边界的情形，平坦 $B$-场以下述方式进入 sigma 模型路径积分。给定映射 $\Phi:\Sigma\to X$，把 $b$ 拉回为 $\Phi^*(b)\in H^2(\Sigma,U(1))=U(1)$，然后在路径积分中乘入因子 $\Phi^*(b)$。于是，在这种情形下，引入平坦 $B$-场的作用是给 $\Sigma$ 到 $X$ 的不同映射分支赋以相位权重。

现在 ${\bf e}_0$ 在映射 $\Phi:\Sigma\to \MH(G_{\rm ad},C)$ 的 sigma 模型中决定一个平坦 $B$-场。事实上，如我们已经解释过的，存在自然的类 $\zeta\in H^2(\MH(G_{\rm ad},C),\CZ)$，它表示把泛 $G_{\rm{ad}}$-丛 ${\EUE}_{\rm {ad}}\to\MH(G_{\rm ad},C)\times C$ 提升为 $\bar G$-丛的障碍。把 $\zeta$ 与 ${\bf e}_0:\CZ\to U(1)$ 复合，我们得到平坦 $B$-场 $b_{{\bf e}_0}={\bf e}_0(\zeta)\in H^2(\MH(G_{\rm{ad}},C),U(1))$。我们断言：${\bf e}_0$ 在映射 $\Sigma\to \MHAD$ 的有效 sigma 模型中的作用，正是按带平坦 $B$-场 $b_{{\bf e}_0}$ 的 sigma 模型所预期的方式给每个映射赋权。

事实上，${\bf e}_0$ 的定义就是路径积分包含相位因子 ${\bf e}_0({\bf a}_0)={\bf e}_0(\xi(E)^{2,0})$。但 $\xi(E)^{2,0}=\Phi^*(\zeta)$；这可由把 \eqref{eq:yrog} 限制到 $C$ 中一点看出。所以 ${\bf e}_0$ 在路径积分中诱导的相位因子是 ${\bf e}_0(\Phi^*(\zeta))=\Phi^*({\bf e}_0(\zeta))=\Phi^*(b_{{\bf e}_0})$。这就证实了我们的断言：${\bf e}_0$ 在 sigma 模型中的作用是生成平坦 $B$-场 $b_{{\bf e}_0}$。

因此 \eqref{eq:niko} 意味着 $S$-对偶把平坦 $G_{\rm{ad}}$-丛的拓扑类 ${\bf m}_0$ 与由 ${\bf e}_0$ 决定的平坦 $B$-场互换。

\begin{yizhu}[物理图像]
按：这里把两个量分得很清楚：${\bf m}_0$ 是离散拓扑数据，它标记 $\MHAD$ 的连通分支（即 $E$ 在 $C$ 上的拓扑类型）；${\bf e}_0$ 是相位数据，它经由 $\zeta$ 给出平坦 $B$-场 $b_{{\bf e}_0}$，从而给不同映射分支加权。式 (7.24) 说的正是 $S$-对偶把这两者互换。参见 Witten 2008 综述关于离散 $\theta$ 角与 $B$-场的讨论。
\end{yizhu}

\par\medskip\emph{膜的引入}

对于几何 Langlands 纲领的应用而言，真正的收获在于理解这一切对膜的含义。

于是我们取 $\Sigma$ 为带边界的黎曼面。$M=\Sigma\times C$ 因此是带边界 $\partial M=\partial\Sigma\times C$ 的四维流形。在 $\partial M$ 上，我们放置某个超对称边界条件。

有效的二维描述是映射 $\Phi:\Sigma\to \MHAD$ 的 sigma 模型，其边界条件对应某个膜。我们想理解 ${\bf e}_0$ 与 ${\bf m}_0$ 在这一描述中的作用。

关于 ${\bf m}_0$ 没有什么新东西可说。无论 $\Sigma$ 是否有边界，它都标记 $\MHAD$ 的连通分支。

${\bf e}_0$ 的作用则更微妙。如我们所见，靶为 $\MHAD$ 的 sigma 模型配备有平坦 $B$-场 ${\bf e}_0(\zeta)$。平坦 $B$-场对膜有非常有趣的影响 \cite{ref102}。没有 $B$-场时，$\MHAD$ 上的膜有 Chan-Paton 丛，它是 $\MHAD$ 上的向量丛（更一般地，是可能支撑在子流形上的层，它定义 $\MHAD$ 的一个 $K$-理论类）。然而，当存在平坦 $B$-场（对应元素 $b\in H^2(\MHAD,U(1))$）时，Chan-Paton 丛变成扭化向量丛（更一般地，是与 $\MHAD$ 的扭化 $K$-理论中某元素相关的扭化层），其扭化由 $b$ 按式 \eqref{eq:hofto} 的意义给出。正是因为这个，我们才引入了扭化向量丛的概念。

所以简而言之，当 ${\bf e}_0\not=0$ 时，膜的 Chan-Paton 丛是扭化向量丛，扭化为 $b_{{\bf e}_0}={\bf e}_0(\zeta)$。所幸，由第~\ref{sec:reviewthis}~节末的分析，我们有大量这样的扭化向量丛可用。若 $\varrho$ 是 $\bar G$ 的任意不可约表示，使得由 $\varrho$ 定义的 $\bar G$ 中心的特征标等于 ${\bf e}_0$，则表示 $\varrho$ 中的泛丛 ${\EUE}_\varrho$ 就是平坦 $B$-场 $b_{{\bf e}_0}$ 的一个扭化向量丛的例子。

其他重要的扭化膜例子可以这样构造：选取子流形 $Y\subset \MHAD$，使得 $b_{{\bf e}_0}$ 限制到 $Y$ 上平凡。此时，支撑在 $Y$ 上的膜的定义与 ${\bf e}_0$ 无关（相差一个并不完全典范的同构）。例如，$Y$ 可以是 $\MHAD$ 中的一点。$b_{{\bf e}_0}$ 限制到一点上当然是平凡的，所以对任意 ${\bf e}_0$，零膜都存在。由这样的零膜，我们可以构成电特征膜，这将在第~\ref{sec:electricbranes}~节解释。稍微更微妙的例子是磁特征膜，即支撑在 Hitchin 纤维化的一条纤维 ${\CMF}$ 上的秩一膜。这种膜的 Chan-Paton 丛应当是平坦\footnote{这里可以把「平坦」理解为与普通线丛一样，即过渡函数是常数。}扭化线丛。文献 \cite{ref21} 证明了平坦 $b$-场 $b_{{\bf e}_0}$ 限制到 ${\CMF}$ 上平凡，因此平坦扭化线丛的空间与 ${\CMF}$ 上普通平坦线丛的空间同构，但并非典范同构。

现在我们可以从 $S$-对偶推出一个关于膜的陈述。对偶变换 $S$ 把 $({\bf e}_0,{\bf m}_0)\to (-{\bf m}_0,{\bf e}_0)$，所以它把「膜所支撑的 $\MH$ 分支的拓扑」与「其 Chan-Paton 丛扭化所依据的平坦 $B$-场 $b_{{\bf e}_0}$」互换。当然，它还把 $G_{\rm ad}$ 与\footnote{所谓 $^L\negthinspace G_{\rm ad}$，我们指群 $^L\negthinspace G$ 的伴随形式。这里描述的陈述最好用双侧的伴随丛来表述，以便容许所有可能的拓扑。我们只在暂时需要时会明确说明，所取的是 $\MH$ 的某个分支还是它的万有覆盖。} $^L\negthinspace G_{{\rm ad}}$ 互换，并且（如我们在本节结尾评注中所解释的）把 $\MH(G_{{\rm ad}})$ 与 $\hat\M_H({}^L\negthinspace G_{{\rm ad}})$ 互换，后者我们定义为 $\MH({}^L\negthinspace G_{{\rm ad}})$ 的万有覆盖。

作为这种对偶的特例，一侧位于 Hitchin 纤维化某条纤维 ${\CMF}$ 中的一点，被映射到另一侧支撑在 Hitchin 纤维化相应纤维 $^L\negthinspace{\CMF}=\Xi({\CMF})$ 上、并配有平坦扭化线丛的秩一膜。${\CMF}$ 是一族复环面之并，由 Higgs 丛的特征类 ${\bf m}_0=\xi(E)$ 标记。一侧 ${\CMF}$ 分支的选取，决定了另一侧的离散电场 ${\bf e}_0$，从而决定了扭化。${\CMF}$ 上一点的选取，决定了 $^L\negthinspace{\CMF}$ 上的平坦扭化线丛（在 $S$-对偶变换下它被映到那里）。当然，${\CMF}$ 与 $^L\negthinspace{\CMF}$ 之间的这种关系是相互的。${\CMF}$ 与 $^L\negthinspace{\CMF}$ 之间的这种扭化对偶，实际上是 Hausel 与 Thaddeus \cite{ref21} 的主要结果之一。

\begin{yizhu}[易错点]
按：$^L\negthinspace G_{\rm ad}$ 指 Langlands 对偶群 $^L\negthinspace G$ 的伴随形式，而不是 $G_{\rm ad}$ 本身的 Langlands 对偶群。原文选择在两侧都用伴随丛来陈述对偶，是为了容纳所有可能的拓扑；否则就得同时选定 $\MH$ 的某个分支或它的万有覆盖。$\hat\M_H$ 是 $\MH$ 的万有覆盖，两者的区别在第 7.2 节末的评注中还会出现。
\end{yizhu}

\par\medskip\emph{${\bf e}_1$ 与 ${\bf m}_1$ 的解释}

最后让我们用二维的语言讨论 ${\bf e}_1$ 与 ${\bf m}_1$ 的解释。先比较规范群为 $\bar G$ 与 $G_{\rm{ad}}$ 的两个理论会有帮助。

$\MH(\bar G)$ 是单连通的。它有一个自然的对称群 ${\bf E}_C=H^1(C,\CZ)$。事实上，${\bf E}_C$ 参数化 $C$ 上的 $\CZ$-丛。一个 $\bar G$ Higgs 丛 $(E,\varphi)$ 可以与一个 $\CZ$-丛作张量积，得到新的 $\bar G$ Higgs 丛。（具体地说，这个操作把 $E$ 绕 $C$ 中非可缩闭圈的 holonomy 乘以 $\bar G$ 中心的元素。）所以这给出靶为 $\MH(\bar G)$ 的 sigma 模型的一个对称群 ${\bf E}_C$。该 sigma 模型的 Hilbert 空间可以按 ${\bf E}_C$ 的特征标分解。

另一方面，$\MHAD$ 没有这样的几何对称性。但它有基本群 ${\bf M}_C=H^1(C,\CZ)$，与 ${\bf E}_C$ 同构。要理解这个基本群从何而来，一个捷径是注意到 $\MHAD$ 的一个分支 $\MHAD_0$，即参数化可提升为 $\bar G$ 的 Higgs 丛的那个分支，恰好是 $\MHAD_0=\MH(\bar G)/{\bf E}_C$。除以 ${\bf E}_C$ 消去了 $\MH(\bar G)$ 的几何对称性，但当然也产生了基本群。\footnote{${\bf E}_C$ 的作用并不自由，所以 $\M_H(G_{ad})$ 有轨形奇点。（它还有来自可约 Higgs 丛的更严重奇点。）正如 sigma 模型中所熟悉的，$\M_H(G_{ad})$ 的基本群必须在轨形意义下理解。或者，也可以依赖四维规范场论，而不必约化到带有奇点的 sigma 模型。}所以 $\pi_1(\MHAD_0)\cong {\bf E}_C$。

实际上，$\MHAD$ 任一连通分支的基本群都相同。$\MHAD$ 的定义是：把规范群为 $G_{\rm {ad}}$ 的 Hitchin 方程的所有解空间，除以 $C$ 上所有 $G_{\rm{ad}}$-值规范变换构成的群 ${\cal G}_{\rm{ad}}(C)$。如果只除以 ${\cal G}_{\rm{ad}}(C)$ 的连通分支，得到的就是万有覆盖 $\hat\M_H(G_{\rm {ad}},C)$。因此 $\MHAD$ 的基本群就是 ${\cal G}_{\rm ad}(C)$ 的分支群，而对 $\MHAD$ 的任一分支，这个群都是 ${\bf E}_C=H^1(C,\CZ)$。

总之，在 $\MHAD$ 上运动的弦有一个离散的守恒缠绕数群 ${\bf M}_C=\pi_1(\MHAD)=H^1(C,\CZ)$。同样，在 $\MH(\bar G)$ 上运动的弦有一个离散的守恒动量群 ${\bf E}_C=H^1(C,\CZ)$。

让我们把对称性 ${\bf E}_C$ 与 ${\bf M}_C$ 同我们在底层规范场论中能看到的东西作比较。在 $M=\Sigma\times C=S^1\times \tilde S^1\times C$ 上的 $G_{\rm{ad}}$ 规范场论中，丛 $E$ 及其他数据决定映射 $\Phi:S^1\times \tilde S^1\to \MHAD$。在分析这一情形的拓扑时，我们把特征类 $\xi(E)$ 展开为 $\xi(E)={\bf a}\oplus {\bf m}$，其中 ${\bf m}$ 是 $\xi(E)$ 在 $\tilde S^1\times C$ 上的限制（更确切地说，是针对 $p\in S^1$ 限制到 $p\times \tilde S^1\times C$）。在低能 sigma 模型中，${\bf a}$ 与 ${\bf m}$ 是 $\Phi:\Sigma\to \MHAD$ 的不变量。特别地，${\bf m}$ 只依赖于 $\Phi$ 在 $\tilde S^1$ 上的限制 $\Phi|_{\tilde S^1}$。

在 \eqref{eq:nimot} 中，我们进一步展开了 ${\bf m}={\bf m}_0\oplus {\bf m}_1$，其中 ${\bf m}_1\in H^1(\tilde S^1,H^1(C,\CZ))$。${\bf m}_1$ 是 $\Phi|_{\tilde S^1}$ 的一个拓扑不变量，当 $\tilde S^1$ 到 $\MHAD$ 的映射为常值时它为零。所以它度量映射 $\Phi|_{\tilde S^1}$ 在 $\pi_1(\MHAD)=H^1(C,\CZ)$ 中的同伦类。

类似地，我们可以交换 $S^1$ 与 $\tilde S^1$ 的角色。$\xi(E)$ 在 $S^1\times C$ 上的限制（即对 $q\in \tilde S^1$ 限制到 $S^1\times q\times C$），在上述记号中是 ${\bf m}_0\oplus {\bf a}_1$。特别地，${\bf a}_1\in H^1(S^1,H^1(C,\CZ))$ 度量 $\Phi$ 在 $S^1$（即「时间」）方向上的缠绕。因此与 ${\bf a}_1$ 对偶的特征标 ${\bf e}_1$ 度量弦的守恒动量。

\begin{yizhu}[物理图像]
按：$\MH(\bar G)$ 上的几何对称群 ${\bf E}_C$ 与 $\MHAD$ 的基本群 ${\bf M}_C$ 都同构于 $H^1(C,\CZ)$，但物理含义不同：前者给出弦的守恒动量，后者给出弦的守恒缠绕数——正如紧致圆上弦的动量与缠绕。$S$-对偶（Langlands 对偶）把这两者互换，这正是它与二维 $T$-对偶（半径 $R\leftrightarrow 1/R$）相关的表现。参见 Frenkel 2009 综述中关于 $S$-对偶与镜像对称的讨论。
\end{yizhu}

\par\medskip\emph{一个例子}

作为这一切的一个重要说明，考虑 Langlands 对偶对 $G_{\rm{ad}}$ 与 $^L\negthinspace\bar G$。在 $G_{\rm{ad}}$ 规范场论中，我们令 ${\bf e}=0$，并考虑靶为 $\MHAD$ 的 sigma 模型。在这个模型中，没有守恒动量，但有关于弦缠绕的对称群 ${\bf M}_C$。

在对偶图景中，规范群为 $^L\negthinspace\bar G$，我们令 ${\bf m}=0$，靶为 $\MH(\bar G)$ 的 sigma 模型没有守恒的弦缠绕，但有一个离散守恒动量的对称群 ${\bf E}_C$。

$S$-对偶、或 Montonen-Olive 对偶、或 Langlands 对偶交换这两幅图景，把 ${\bf E}_C({}^L\negthinspace\bar G)$ 与 ${\bf M}_C(G_{\rm{ad}})$ 互换。这个对偶交换弦的离散守恒动量与缠绕这一事实，是它与二维 $T$-对偶之关系的一个侧面。

\par\medskip\emph{结尾评注}

更一般地，我们可以随意指定 ${\bf m}_0$ 与 ${\bf e}_0$。然后我们考虑按 ${\bf m}_0$ 作拓扑标记的 $\MHAD$ 的某个分支上的膜；这些膜的 Chan-Paton 丛被 ${\bf e}_0$ 扭化。

我们仍有两条量子化途径。如果除以所有 $G_{\rm {ad}}$-值规范变换，则靶空间是按 ${\bf m}_0$ 标记的 $\MH(G_{\rm ad},C)$ 的分支。此时有一个分类弦缠绕数的有限群 ${\bf M}_1$，但没有几何对称群 ${\bf E}_1$。如果只除以连通规范变换，则靶是万有覆盖 $\hat\M_H(G_{\rm ad},C)$。此时有几何对称群 ${\bf E}_1$，但没有弦缠绕数群 ${\bf M}_1$。$S$-对偶交换 ${\bf E}_1$ 与 ${\bf M}_1$，所以（除了交换 $G$ 与 $^L\negthinspace G$ 之外）它还交换两种量子化方式。

特别地，当我们把 $S$-对偶作用于给定 ${\bf e}_0$ 与 ${\bf m}_0$ 的膜时，除了交换两个伴随群 $G_{\rm ad}$ 与 $^L\negthinspace G_{\rm ad}$ 之外，还必须交换两种量子化方式。

\begin{yizhu}[含糊点拨]
按：本节的论证依赖一个容易被忽略的环节：$\MHAD$ 上「除以全部规范变换」与「只除以连通规范变换」给出不同的靶空间（前者是某个分支，后者是其万有覆盖），二者的对称性也互补——一个只有缠绕数群，另一个只有几何对称群。$S$-对偶必须同时交换这两者，这也是原文结尾评注的要点。
\end{yizhu}
